% TOP_PROBLEM: {"id": 3231, "title": "Partial List Coloring", "category_id": 3, "category": {"id": 3, "name": "graph_theory", "display_name": "Graph Theory", "description": "Problems involving graphs, networks, and their properties.", "slug": "graph-theory", "order_index": 3, "created_at": "2026-07-31T15:26:25.671Z"}, "status": "open", "problem_number": "OPG-771", "set_id": 12}
% FIRST_POSTED: 2026-10-01
% ENTRY_KIND: statement_only
% UnsolvedMath ID 3231; source catalogue code OPG-771
% !TeX program = lualatex
% Definition and sources only; this file is not an LLM solution attempt.
\documentclass[11pt,a4paper]{article}
\usepackage[margin=25mm]{geometry}
\usepackage{fontspec}
\setmainfont{lmroman10-regular.otf}[BoldFont=lmroman10-bold.otf,
 ItalicFont=lmroman10-italic.otf,BoldItalicFont=lmroman10-bolditalic.otf]
\usepackage{amsmath,amssymb,mathtools}
\usepackage{unicode-math}
\setmathfont{latinmodern-math.otf}
\AtBeginDocument{\let\setminus\smallsetminus}
\usepackage{xurl,hyperref}
\hypersetup{colorlinks=true,urlcolor=blue}
\setcounter{secnumdepth}{0}
\setlength{\parindent}{0pt}
\setlength{\parskip}{5pt}
\setlength{\emergencystretch}{3em}
\begin{document}
\section*{Partial List Coloring}\label{3231}
{\small UnsolvedMath ID 3231 · OPG-771 · Graph Theory · OpenGarden}
\subsection{Definitions and mathematical statement}
Let $G=(V,E)$ be a nonempty
finite simple graph with
$n=|V|$. A list assignment
gives each vertex $v$ a
finite set $L(v)$ of allowed
colours. An $L$-colouring
is a function $c$ with
$c(v)\in L(v)$ at every
vertex and $c(u)\neq c(v)$
on every edge. The list
chromatic number $\chi_\ell(G)$
is the least integer $q$
such that every list assignment
with at least $q$ colours
at each vertex permits a
proper $L$-colouring.

Put $q=\chi_\ell(G)$, so
$q\geq1$, and let $t$
be an integer with
$0\leq t\leq q$. A partial
proper list colouring consists
of a subset $U\subseteq V$
and a proper list colouring
of the induced graph $G[U]$
using the original lists.
Vertices outside $U$ are
left uncoloured, so they
impose no colour choice.

The conjecture asserts that
for every assignment of
exactly $t$ colours to each
vertex there is such a
partial colouring with
\[
 |U|\geq\frac{tn}{q},
 \quad\text{equivalently}\quad
 |U|\geq\left\lceil\frac{tn}{q}\right\rceil.
\]
The lists may differ arbitrarily,
and $U$ may be selected
after seeing them. The cases
$t=0$ and $t=q$ respectively
allow the empty colouring
and require a full colouring
already guaranteed by the
definition. Intermediate list
sizes request the stated
linear proportion uniformly
over all graphs and all
lists, not merely in
expectation.
\subsection{Short English statement}
With lists of t colours, can every graph colour at least the fraction t divided by its list chromatic number of its vertices?
\subsection{Sources}
\begin{itemize}
\item[S1] ULAM.AI and the UnsolvedMath Contributors, supplied problems.json, record 3231 (OPG-771), OpenGarden collection. Catalogue identity and source transcription; CC BY 4.0.
\url{https://huggingface.co/datasets/ulamai/UnsolvedMath}.
\item[S2] Open Problem Garden, Partial List Coloring. Original problem and discussion; the mathematical formulation below is expanded from the supplied source record.
\url{http://www.openproblemgarden.org/op/partial_list_coloring}.
\end{itemize}
\end{document}
